Koszul complexes. The Koszul complex associated with the sequence @@M0@@ is isomorphic to the tensor-product complex @@M1@@. Since it is a resolution of @@M2@@ by free @@M3@@-Modules of finite type, the @@M4@@ are the homology sheaves of the complex @@M5@@. We can regard it as the tensor product of the complexes @@M6@@ and @@M7@@, and compute its homology using the spectral sequences associated with a bicomplex (EGA @@M8@@ (11.3.5)), which abut to the homology of the associated simple complex.

Consider the spectral sequence such that @@M9@@, in the notation of EGA @@M10@@ (11.3.5), with the first index referring to the complex @@M11@@ and the second to the complex @@M12@@. Now the complex @@M13@@ is the Koszul complex over @@M14@@ associated with the sequence of sections @@M15@@, where the @@M16@@ are the images of the @@M17@@ in @@M18@@, and this complex is a resolution of @@M19@@ over @@M20@@ since the immersions @@M21@@, and @@M22@@ are regular (cf. the proof of 1.7 ii). Since @@M23@@ is free over @@M24@@, the complex @@M25@@ is a resolution of @@M26@@. Consequently the spectral sequence degenerates, and the @@M27@@ are the homology sheaves of the complex @@M28@@. The differential of @@M29@@ is expressed in terms of the @@M30@@, which annihilate @@M31@@; consequently the differential of @@M32@@ is zero, and we therefore obtain the isomorphisms:
\sgaStableTarget{sga6.e07.eq.2-5-2}@@M33@@

In particular, for @@M34@@ we find:
@@M35@@
whence for every @@M36@@:
\sgaStableTarget{sga6.e07.eq.2-5-3}@@M37@@
